\section{Introduction}
\label{sec:intro}

Terrestrial life evolved over nearly four billion years within the continuous, dynamic presence of Earth's geomagnetic field (GMF), which provides an omnipresent flux density of $\sim\SI{25}{\micro\tesla}$ to \SI{65}{\micro\tesla} at the surface. The GMF serves a dual evolutionary role: on a planetary scale, it defers the solar wind and deflects destructive solar energetic particles; on a biochemical and physiological scale, it acts as an environmental stimulus influencing fundamental biological processes ranging from cryptochrome-mediated light sensing and radical pair spin dynamics to circadian rhythm entrainment, root gravitropism, and cellular redox homeostasis \citep{Maffei2014, Binhi2017, Hore2016}.

In sharp contrast, the magnetic environment of Mars underwent a catastrophic divergence roughly 4.1 to 3.9 billion years ago during the late Noachian era, when its internal core dynamo ceased operation \citep{Acuna1999, Connerney2005}. Today, Mars lacks an active global dipole field. However, unlike the Moon, Mars preserves intense, large-scale remnant crustal magnetization embedded within its ancient basaltic and Noachian crust \citep{Connerney2005, Langlais2019}. These crustal magnetic anomalies exhibit extreme hemispheric dichotomy: the ancient southern cratered highlands (particularly Terra Cimmeria and Terra Sirenum) harbor coherent magnetic lineations with field strengths exceeding \SI{200}{\nano\tesla} to \SI{300}{\nano\tesla} at \SI{400}{\kilo\meter} orbital altitudes and $> \SI{1500}{\nano\tesla}$ at the surface \citep{Acuna1999, Langlais2019}. In contrast, the younger northern lowlands (Vastitas Borealis, Utopia Planitia, Arcadia Planitia) and large impact basins (Hellas, Argyre, Isidis) underwent extensive shock, thermal, and hydrothermal demagnetization, leaving vast regions with near-null background crustal fields ($< \SI{5}{\nano\tesla}$ to \SI{20}{\nano\tesla}) \citep{Langlais2019, Morschhauser2014}.

From a space biology and physiological perspective, any planetary environment with magnetic intensities substantially below Earth's baseline ($< \SI{5}{\micro\tesla}$) constitutes a hypomagnetic field (HMF). While the interplanetary cruise phase exposes biological systems to an isotropic interplanetary magnetic field (IMF) of $\sim\SI{5}{\nano\tesla}$, surface operations on Mars will experience extreme geographic heterogeneity depending on landing site selection. Experimental HMF studies on Earth demonstrate that prolonged magnetic deprivation induces significant phenotypic and molecular stress, including delayed flowering time in model plants, impaired polar auxin transport in roots, disrupted iron uptake, altered gut microbiome taxonomic ratios, and exacerbated musculoskeletal bone demineralization in mammalian models \citep{Agliassa2018, Narayana2018, Narayana2021, Zhang2017, Xu2012, Mo2014}.

As international space programs and commercial initiatives actively prepare for human exploration of Mars—focusing heavily on ice-rich northern candidate sites such as Arcadia Planitia and Deuteronilus Mensae—the physiological implications of Martian magnetic field heterogeneity remain largely uncharacterized. This study addresses this critical gap by performing a global spatial analysis of the Martian crustal magnetic field, quantitatively extracting the magnetic environment at 13 active, historical, and candidate human landing sites, and synthesizing these geophysical constraints with an exhaustive review of terrestrial magnetobiology literature to formulate an actionable biological risk framework for future surface operations.
